\section{The normalized indicators, covariance, and canonical finite jet field}
\label{sec:canonical-analytic}

\status{The claims in this section reconstruct the actual analytic and finite
Fourier identities used in \texttt{finAlg\_step}, and then join them to the
canonical local-index construction on the same data. The finite operator
claims CA1--CA4 do not assume a theta vacuum. The actual vacuum and its
algebraic translates are supplied by VI. The normalization in covariance
is proved, including its determinant character and its inert-prime sign;
it is not deduced from theta-line uniqueness alone. The final construction
retains the literal frame-inclusion and raw \texttt{PreX} inputs, and does
not construct a later geometric stage. All statements are human
reconstructions from the pinned source, without a Lean build or a
kernel-verification claim.}

\paragraph{Conventions.}
Fix the source datum $D$, its CM field $E$, imaginary quadratic subfield
$K=\mathbb Q(\sqrt{-d})$, induced CM type $\Phi$, and a curve ideal
$\mathfrak c$ with discriminant normalization $\delta$. Write
$V=\prod_{\tau\in\Phi}\operatorname{Mat}_{b\times a}(\mathbb C)$ and
$H$ for the source Hermitian form, linear in its first argument:
\[
 H(z,z')=\frac1{\operatorname{area}(\delta\mathfrak c)}
       \sum_{\tau,i,j}z_{\tau ij}\overline{z'_{\tau ij}},\qquad
 E_H(z,z')=\operatorname{Im}H(z,z'),\qquad
 Q(z)=\frac{\sqrt d}{2}H(z,z).
 \tag{CA.1}
\]
The last quantity is real. Rational points mean the image
$X=\delta\operatorname{Emb}_{\Phi}\operatorname{Mat}_{b\times a}(E)$,
not arbitrary complex points. For self-dual Hermitian lattices $L,L'$,
write $\Lambda(L,L')$ for the source period lattice \texttt{La}; under
this rational identification it is
$\mathfrak c\,\mathfrak A\operatorname{Hom}_{\mathcal O_E}(L,L')$,
where $\mathfrak A$ is the source trace-dual ideal. Every assertion about
this identification uses that same ideal and that same $\delta$.
For additive lattices $A,B$, set $[B:A\cap B]=n(A,B)$; this is the
source relative-index convention.

\paragraph{CA1: the exact rational Heisenberg operators.}
Define
\[
 (U_x f)(z)=e^{-\pi H(z,x)-\pi H(x,x)/2}f(z+x),\qquad
 \alpha(x)=e^{\pi iQ(x)},\qquad W_x=\alpha(x)U_x.
 \tag{CA.2}
\]
Expanding $H(z+x,y)$ and $H(x+y,x+y)$, using
$H(y,x)=\overline{H(x,y)}$, gives
\[
 U_xU_y=e^{-\pi iE_H(x,y)}U_{x+y},\qquad
 W_xW_y=e^{-2\pi i c(x,y)}W_{x+y},\qquad
 c(x,y)=\tfrac12\bigl(\sqrt d\,\operatorname{Re}H(x,y)
                         +E_H(x,y)\bigr).
 \tag{CA.3}
\]
These formulas hold on every function, so they also hold termwise in
finite sums and under differentiation of entire functions. In particular
$W_xW_y=e^{-2\pi iE_H(x,y)}W_yW_x$.

The integer $w=(1+\sqrt{-d})/2\in\mathcal O_K\subset\mathcal O_E$
has $w+\bar w=1$; all embeddings in $\Phi$ send it to
$(1+i\sqrt d)/2$. Thus the last cocycle in (CA.3) is
$\gamma(wx,y)$, where
\[
 \gamma(x,y)=e^{-2\pi iE_H(x,y)}.
 \tag{CA.4}
\]
It is alternating and biadditive, and
$\gamma(rx,y)=\gamma(x,\bar r y)$ for $r\in\mathcal O_E$.
The existence and integrality of $w$ are the actual quadratic-order
facts in the host, including the prime $2$; no division by $2$ is
performed in a residue field.

For rational $x,y$, both $E_H(x,y)$ and $\sqrt d\,
\operatorname{Re}H(x,y)$ are rational. Indeed (CA.1) becomes the
trace of an element of $E$ to the induced quadratic field, divided
by the rational ideal norm and by $\sqrt d/2$; its imaginary and
$\sqrt d$ times its real part are rational. Consequently the phases
in (CA.2)--(CA.4) are roots of unity. All their algebraic conjugates
have modulus one. This is the exact rationality used below, not
algebraicity of an arbitrary complex Gaussian factor.

\paragraph{CA2: a finite presentation of the normalized indicator.}
Use the literal \texttt{ambOf} datum obtained from a seam $s$, with
\texttt{Fields.Ok}, \texttt{Seam.Ok}, the above curve data, self-dual
$L,L'$, an arbitrary finite prime set $S$, and an integer $c>0$.
Let $\Lambda=\Lambda(L,L')$, $A_c=\texttt{AC}$,
$B_c=\texttt{BC}$, and choose the actual representative sets
\[
 R_A\subset\Lambda\cap A_c\quad\bmod B_c,
 \qquad R_B\subset B_c\quad\bmod\Lambda.
 \tag{CA.5}
\]
Here ``modulo'' means modulo the intersection when neither lattice
contains the other. Let $s_x$ be the source's chosen rational lift of
$x$, let $v=\texttt{volC}$, and put $k=\texttt{chiF}(L,L')$.
Writing $G=\texttt{thMoved}(L,L')$ in host coordinates, the actual
definition of the indicator expands to
\[
 k^{-1}\operatorname{thInd}_{S,L,L',c,x}
       =\sum_{y\in R_B}r_yU_{y-s_x}G,
 \tag{CA.6}
\]
where
\[
\begin{split}
 r_y={}&(k^{-1}k)\,\frac{v}{(\sqrt v)^2}\,
 \left(\frac1{|R_A|}\sum_{t\in R_A}
       {\bf1}_{B_c}(t)e^{2\pi iE_H(t,y)}\right)\\
 &\hspace{12mm}\cdot\alpha(s_x)\alpha(y)
                      e^{\pi iE_H(s_x,y)}.
\end{split}
 \tag{CA.7}
\]
If a degenerate definition gives $v\le0$, the two inverse square
roots in the source make their combined factor zero; in the present
full-lattice situation $v>0$ and the displayed factor is one. Likewise
$k^{-1}k$ is zero or one, with no algebraicity assumption on $k$.
For an actual character satisfying \texttt{IsChi}, it is one.

To check the normalization, expand the source finite Fourier
transform \texttt{hatS}, then its two factors $(\sqrt v)^{-1}$.
The translation identity
\[
 (U_yG)(z-s_x)e^{-\pi H(z,-s_x)-\pi H(-s_x,-s_x)/2}
       =e^{\pi iE_H(s_x,y)}(U_{y-s_x}G)(z)
\]
is CA1. It gives precisely (CA.7); in particular the inverse character
has canceled before any coefficient is called algebraic.

Every $r_y$ is algebraic and every algebraic embedding sends it to a
number of modulus at most one: the inner expression is an average
of zeros and roots of unity, and the remaining nonzero factors have
modulus one under every embedding. Every $y-s_x$ is rational with
respect to $\Lambda$.

The representative sets exist and satisfy the concrete estimate
\[
 |R_B|\le c^{[E:\mathbb Q]ab}.
 \tag{CA.8}
\]
For completeness, the local exponents defining either half lattice
lie between $-1$ and $1$. Hence $c\Lambda\subset B_c$ and
$cB_c\subset\Lambda$. The map $y\mapsto cy$ injects
$B_c/(B_c\cap\Lambda)$ into $\Lambda/c\Lambda$. A full rational
lattice has a $\mathbb Z$-basis of size $[E:\mathbb Q]ab$, so the
latter quotient has $c^{[E:\mathbb Q]ab}$ elements. The same
inclusions prove finiteness for $R_A$. This argument also covers
$a=0$ or $b=0$, when the quotients have one element.

Thus (CA.6) is the literal source predicate \texttt{Pres}, with the
sum of the absolute values of the conjugated coefficients bounded
by (CA.8). It holds for every $S$ here; the stronger no-conjugate-pair
hypothesis needed for \texttt{LaHyp} and covariance is not inserted
into this finite-presentation assertion.

\paragraph{CA3: the finite Hermitian norm Gauss sum, including characteristic two.}
Let $k/k_0$ be a finite field with nontrivial involution $\sigma$,
$|k_0|=q$, and let $\eta:k_0\to\mathbb C^\times$ be a nontrivial
additive character. Then
\[
 \sum_{t\in k}\eta(t\sigma(t))=-q=-\sqrt{|k|}.
 \tag{CA.9}
\]
Indeed $k^\times$ is cyclic and the norm $t\mapsto t^{q+1}$ maps
onto $k_0^\times$, with every nonzero fiber of size $q+1$.
Translation by a point at which $\eta\ne1$ shows
$\sum_{u\in k_0}\eta(u)=0$. The left side is consequently
$1+(q+1)(-1)=-q$.

The precise form needed for the Heisenberg cocycle is as follows.
Let $\chi:k\to\mathbb C^\times$ be nontrivial, additive, and
$\sigma$-invariant, and suppose $w+\sigma(w)=1$.
Then $\chi=\eta\circ\operatorname{Tr}_{k/k_0}$ for a nontrivial
$\eta$, and
\[
 \sum_{t\in k}\chi(wt\sigma(t))=-\sqrt{|k|}.
 \tag{CA.10}
\]
The kernel of the trace is the image of $\sigma-1$, on which
$\chi$ is trivial. The trace is onto: $X+X^q$ is a nonzero polynomial
of degree $q<q^2$, even in characteristic two. This proves the
factorization and shows its character is nontrivial. Since
$t\sigma(t)\in k_0$ and $\operatorname{Tr}(w)=1$, (CA.10) is
(CA.9). There is no missing dyadic exception.

\paragraph{CA4: finite intertwiners with their correct signs compose exactly.}
Let $R=\mathcal O_E$ act on an additive group $X$, let $W_x$ be
linear operators satisfying (CA.3) in the form
$W_xW_y=\gamma(wx,y)W_{x+y}$, and let $A,B,C$ be mutually
commensurable self-dual $R$-submodules for $\gamma$. On an $A$-fixed
vector define
\[
 S_{A,B}v=\sum_{x\in B/(A\cap B)}W_xv,\qquad
 J_{A,B}v=\lambda(n(A,B))^{-1}S_{A,B}v.
 \tag{CA.11}
\]
The representatives do not matter: changing $x$ by an element of
$A\cap B$ has phase one. The vector $S_{A,B}v$ is $B$-fixed.

Choose a multiplicative $\lambda:\mathbb N_{>0}\to\mathbb C^\times$
such that, at every residue field $R/P$,
\[
 \lambda(|R/P|)^2=|R/P|,
 \qquad \bar P=P\Longrightarrow
              \lambda(|R/P|)=-\sqrt{|R/P|}.
 \tag{CA.12}
\]
The involution on a fixed residue field must be nontrivial. These
conditions are exactly \texttt{GoodNorm}. For the actual host they
hold: $E/\mathbb Q$ is abelian and $E/E^+$ is unramified. Over a
rational prime $p$ all primes have the same residue degree $f_p$,
and conjugation fixes either all of them or none, since it commutes
with every Galois automorphism. Set
\[
 \lambda(p)=\begin{cases}
 \sqrt p\,e^{\pi i/f_p},&\text{conjugation fixes the primes over }p,\\
 \sqrt p,&\text{otherwise},
 \end{cases}
\]
and extend multiplicatively. The fixed residue involution is
nontrivial by the actual unramified CM condition, as checked in
DR1/AH. This proves (CA.12) without choosing inconsistent square
roots at different primes.

With this normalization,
\[
 J_{B,C}J_{A,B}v=J_{A,C}v.
 \tag{CA.13}
\]
Here is the finite proof. Self-duality gives perfect pairings between
the two finite quotients by each intersection. Averaging a nontrivial
character gives zero, so
$S_{B,A}S_{A,B}v=n(A,B)v$. If $A,B$ are adjacent, write
$B/(A\cap B)=R/P$ and $q_P=|R/P|$. There are three cases.
If $B=(A\cap B)+(B\cap C)$, direct representatives give
$n(A,C)=q_Pn(B,C)$ and $S_{B,C}S_{A,B}=S_{A,C}$.
If $A=(A\cap B)+(A\cap C)$, the reverse calculation gives
$n(B,C)=q_Pn(A,C)$ and $S_{B,C}S_{A,B}=q_PS_{A,C}$.
The factors in (CA.12) cancel in both cases.

Otherwise $A\cap C=B\cap C$. Let $x$ generate
$B/(A\cap B)$. The perfect pairing supplies $x=a+c_0$ with
$a\in A,c_0\in C$. Put $\chi(t)=\gamma(a,tc_0)$. It descends to
a nontrivial additive character on $R/P$; isotropy of $B$ gives
$\chi(\bar t)=\chi(t)$. Nontriviality forces $\bar P=P$:
if some $r\in P$ had $\bar r\notin P$, the inverse of $\bar r$
modulo $P$ and $\chi(\bar t)=\chi(t)$ would make $\chi$ trivial.
Expanding $W_{tx}v$ by the cocycle, then using the $C$-invariance of
$S_{A,C}v$, gives
\[
 S_{B,C}S_{A,B}v
   =\left(\sum_{t\in R/P}\chi(wt\bar t)\right)S_{A,C}v
   =\lambda(q_P)S_{A,C}v.
\]
Here $n(B,C)=n(A,C)$, so (CA.13) again follows. This is the case
where replacing every signed square root by a positive square root
would lose a sign.

Finally reduce general $A,B$ to adjacency. Choose a simple nonzero
$R$-submodule $Rx$ of $B/(A\cap B)$ and set
$A_1=(A\cap(Rx)^\perp)+Rx$. The perfect pairing proves
$A_1=A_1^\perp$; its quotient over $A\cap A_1$ is $R/P$ and
$A_1=(A_1\cap A)+(A_1\cap B)$. Therefore
$n(A,B)=|R/P|n(A_1,B)$, which strictly decreases the positive
integer index. Induction, the adjacent calculation, and the
intermediate-lattice sum formula establish (CA.13) in finitely many
steps. The argument uses finite modules; it makes no assumption
that a global ideal or lattice is free.

\paragraph{CA5: even rectangular ranks identify the signed and actual intertwiners.}
Use VI2's full, self-dual period lattices in the rational space $X$.
The action of $R=\mathcal O_E$ is entrywise through the embeddings;
CA1 gives the adjoint identity and the integral element $w$ needed
by CA4. On a period lattice $A$, $Q(x)$ is integral and
$\alpha(x)=(-1)^{Q(x)}=\texttt{al}(x)$. Thus its theta functions
are exactly the entire $A$-fixed vectors for $W$, and the source
operator is
\[
 I_{A,B}=n(A,B)^{-1/2}S_{A,B}.
 \tag{CA.14}
\]

Suppose $a,b$ are even. If only $L$ changes in $\Lambda(L,L')$,
then at a prime $P$ the length of the quotient by the intersection
is a multiple of $b$. To verify this, choose an
$\mathcal O_{E,P}$-basis of $L'_P$ and diagonalize the two source
lattices over that discrete valuation ring. The Hom lattice is a
direct sum of $b$ copies of the corresponding source dual lattice.
Tensoring by the common locally principal ideal
$\mathfrak c\mathfrak A$ does not alter its relative index.
If only $L'$ changes, the analogous length is a multiple of $a$.
These are local bases; no global basis is required.

The only nonpositive square-root choices in CA4 occur over the
conjugation-fixed primes. Their signs are therefore
$(-1)^{b\ell_P}$ or $(-1)^{a\ell_P}$, both one. At every other
prime the normalization is already positive. Hence
\[
 J_{\Lambda(L,L'),\Lambda(M,L')}=I_{\Lambda(L,L'),\Lambda(M,L')},
 \qquad
 J_{\Lambda(L,L'),\Lambda(L,M')}=I_{\Lambda(L,L'),\Lambda(L,M')}.
 \tag{CA.15}
\]
In particular, if $A_0=\Lambda(\mathcal O_E^a,\mathcal O_E^b)$
and $\theta_0=\texttt{thZero}$, then the actual definition satisfies
\[
 \texttt{thMoved}(L,L')
   =I_{\Lambda(L,\mathcal O_E^b),\Lambda(L,L')}
      I_{A_0,\Lambda(L,\mathcal O_E^b)}\theta_0
   =J_{A_0,\Lambda(L,L')}\theta_0.
 \tag{CA.16}
\]
CA4 now proves both one-variable transport laws and independence of
the order of these two moves. The equality does not claim that
every positive-normalized intertwiner between arbitrary lattices
composes without a sign.

Rational unitary changes of coordinates commute with these $J$
operators in their natural form. If $T$ preserves $H$, $Q$, and
$X$, then it carries representatives bijectively, preserves the
index, and
$P_TW_x=W_{Tx}P_T$, where $P_Tf(z)=f(T^{-1}z)$. Therefore
\[
 P_TJ_{A,B}=J_{TA,TB}P_T.
 \tag{CA.17}
\]
This is a direct finite-sum identity, independent of a separately
constructed adelic Weil representation.

\paragraph{CA6: computing the scalar by a Gaussian trace.}
We give the scalar calculation used for the actual rational action.
Let $A$ be a full self-dual period lattice, let $u$ be unitary for
$H$, put $a_0=u^{-1}$, and suppose $\det_{\mathbb C}(1-a_0)\ne0$.
Assume $A$ and $uA$ are commensurable and choose $m>0$ such that
\[
 mA\subset(1-a_0)(A\cap a_0^{-1}A).
 \tag{CA.18}
\]
Rationality of $u$ supplies such an $m$ by clearing denominators.
Put $C=\tfrac12(1+a_0)(1-a_0)^{-1}$. For nonzero
$\theta\in\operatorname{Th}(A)$, theta-line uniqueness and the
positive-Iop isometry (VI5) give
\[
 P_u\theta=c_u I_{A,uA}\theta,\qquad |c_u|=1.
\]
The additional information needed to determine its phase is
\[
 c_u=r\,\det_{\mathbb C}(1-a_0)^{-1}
       \sum_{x\in A/mA}\alpha(x)e^{-\pi H(Cx,x)},
 \qquad r>0.
 \tag{CA.19}
\]

Here is a derivation, including the analytic step. Define the kernel
\[
 K_A(z,t)=\sum_{x\in A}\alpha(x)
  \exp\!\left(\pi H(t+x,z)-\pi H(t,x)-\frac\pi2H(x,x)\right).
 \tag{CA.20}
\]
Positive definiteness bounds each summand, uniformly for $z,t$ in
a compact set, by $C_1e^{-C_2\|x\|^2+C_3\|x\|}$ with $C_2>0$.
The number of lattice points in a radius-$R$ ball is $O((1+R)^{2\dim V})$,
by disjoint small balls around lattice points. The bound is summable,
also after every fixed derivative. Thus the kernel is entire in $t$,
and shifting $x$ proves that $K_A(z,\cdot)$ lies in the theta line.
The scalar equation for $\theta$ therefore holds for this kernel.

Take a fundamental domain $F$ for
$A_m=(1-a_0)^{-1}(mA)\subset A\cap a_0^{-1}A$.
For fixed $z$, the kernel equation is
$K_A(z,a_0t)=c_u n(A,uA)^{-1/2}
\sum_y\alpha(y)U_y(K_A(z,\cdot))(t)$, with $y$ representing
$uA/(A\cap uA)$. Set $t=z$ and multiply by
$e^{-\pi H(z,z)}$. For each $y$, expansion of (CA.20) gives
\[
 U_y(K_A(z,\cdot))(z)e^{-\pi H(z,z)}
 =\sum_{x\in A}b(x,y)e^{2\pi iE_H(x+y,z)},
\]
where
$b(x,y)=\alpha(x)\exp(-\pi H(x,y)-\pi H(x,x)/2
-\pi H(y,y)/2+2\pi iE_H(x,y))$.
Each character on the right is periodic modulo $A_m$, because
$x\in A$, $y\in uA$, and $A_m\subset A\cap uA$.
Its integral on $F$ vanishes unless $x+y=0$: a nonzero real
linear functional $E_H(x+y,\cdot)$ gives a nontrivial torus
character, whose Haar integral is zero by translation.
If $y\notin A$ there is no such $x$. If $y\in A$, the single
remaining coefficient is $b(-y,y)=\alpha(-y)$; multiplying by
the outer $\alpha(y)$ gives one. Thus only the zero class of
$uA/(A\cap uA)$ remains, with integral $\operatorname{vol}(F)$.
The finite-Iop side is exactly
$c_u n(A,uA)^{-1/2}\operatorname{vol}(F)$, whereas the other
side is the integral of $K_A(z,a_0z)e^{-\pi H(z,z)}$.

On the other side write each $x\in A$ uniquely as
$s-(1-a_0)y$, with $s\in A/mA$ and $y\in A_m$.
The semicharacter identity converts this shift into $z\mapsto z+y$,
so the integral unfolds to a finite sum of integrals over $V$.
Set $B=(1-a_0)^{-1}$. The integral for $s$ is
\[
 \frac{\kappa}{\det_{\mathbb C}(1-a_0)}
       \alpha(s)e^{-\pi i\operatorname{Im}H(Bs,s)},\qquad \kappa>0.
 \tag{CA.21}
\]
To check (CA.21), diagonalize the unitary $a_0$ in an $H$-orthonormal
basis. Each eigenvalue $\zeta\ne1$ has $\operatorname{Re}(1-\zeta)>0$.
The elementary complex Gaussian integral with linear terms is
\[
 \int_{\mathbb C}e^{-\pi h(1-\zeta)|z|^2+\pi(pz+q\bar z)}\,dz
   =\frac1{h(1-\zeta)}
      e^{\pi pq/(h(1-\zeta))}.
\]
It follows from the two real Gaussian integrals, first for real
positive coefficients and then by holomorphic continuation in the
right half-plane; the dominating Gaussian on compact parameter
sets justifies that continuation and differentiation. Multiplying
the one-coordinate formulas and substituting
$(1-a_0)Bs=s$ leaves precisely the phase in (CA.21).
Equivalently, the remaining exponent identity is
$-H(s,s)/2-H(Bs,a_0^{-1}s)=-i\operatorname{Im}H(Bs,s)$.
Since $C=B-\tfrac12$ and $C$ is skew-adjoint,
$H(Cs,s)=i\operatorname{Im}H(Bs,s)$. Dividing by the first
positive volume factor proves (CA.19). All unfolding interchanges
are justified by the Gaussian bound above and the finite number
of cosets.

\paragraph{CA7: the finite Gauss factor is nonnegative in the even-rank action.}
Now $A=A_0$ is the standard rectangular period lattice and $u$
acts on just one matrix factor. The other factor has even dimension.
Let $P$ project onto its first half and let $S$ interchange the two
halves. They preserve $A$, commute with $a_0,C$, and satisfy
$P^2=P$, $S^2=1$, $SP=(1-P)S$. Both halves are $H$-orthogonal.
For
\[
 \varphi(x)=\alpha(x)e^{-\pi H(Cx,x)}
           =\exp\!\left(\pi i\left(Q(x)-\operatorname{Im}H(Cx,x)\right)\right),
 \tag{CA.22}
\]
the semicharacter identity, (CA.18), and integrality of $E_H$ on
$A$ show $\varphi(x+mz)=\varphi(x)$. These same identities give
$\varphi(x)=\varphi(Px)\varphi((1-P)x)$ and
$\varphi(Sx)=\varphi(x)$.

There is an integral $\ell\in\mathcal O_E$ with
\[
 \ell\bar\ell=-1+4m b_0,\qquad b_0\in\mathcal O_{E^+}.
 \tag{CA.23}
\]
At a split local pair choose components $1,-1$; at a nonsplit
unramified place, DR1's norm surjectivity gives norm $-1$ modulo
$4m$. The Chinese remainder theorem combines these finitely many
choices into $\ell$. The quotient $b_0$ is integral and fixed by
conjugation. Thus (CA.23) includes all primes dividing $2m$.

Let $L_0$ be entrywise multiplication by $\ell$, and $B_0$ by $b_0$.
They preserve $A$ and commute with $P,a_0,C$; moreover
\[
 H(L_0x,L_0y)=-H(x,y)+4mH(B_0x,y).
\]
It follows that $L_0$ induces a bijection of $A/mA$: if
$L_0x\in mA$, the displayed pairing identity makes
$E_H(x,y)$ divisible by $m$ for every $y\in A$, and self-duality
gives $x\in mA$. A direct exponent check gives
$\varphi(L_0x)=\overline{\varphi(x)}$. In detail, write
$mx=y-a_0y$ as in (CA.18), so $2mCx=y+a_0y\in A$.
The difference between the two exponents, divided by $\pi i$, is
\[
 4m\left(\tfrac{\sqrt d}{2}\operatorname{Re}H(B_0x,x)
                 -\operatorname{Im}H(B_0Cx,x)\right)\in2\mathbb Z.
\]
The first term is $2m$ times the integer
$Q(B_0x+x)-Q(B_0x)-Q(x)$; the second is twice the integral
pairing of $B_0(y+a_0y)$ with $x$.

The finite quotient $A/mA$ is the direct sum of the image and
kernel of $P$. Its two Gauss sums are equal by $S$. The first is
real because the bijection $L_0$ conjugates its summands and
preserves that half. Therefore their product, the sum in (CA.19),
is the square of a real number and is nonnegative. It cannot be
zero, since $|c_u|=1$. Equation (CA.19) consequently gives
\[
 c_u=\frac{|\det_{\mathbb C}(1-u^{-1})|}
               {\det_{\mathbb C}(1-u^{-1})}.
 \tag{CA.24}
\]
This positivity calculation is the step that fixes the phase; a
unit-modulus assertion alone would not give (CA.24).

\paragraph{CA8: the exact determinant powers for every rational unitary matrix.}
For $\gamma\in U(a,E)$ and $\gamma'\in U(b,E)$ put
\[
 (T_{\gamma,\gamma'}z)_\tau
       =\tau(\gamma')z_\tau\tau(\gamma)^{-1},\qquad
 P_{\gamma,\gamma'}f(z)=f(T_{\gamma,\gamma'}^{-1}z).
 \tag{CA.25}
\]
This is exactly the source \texttt{pull}: the inverse on the left
is the conjugate transpose, whereas the factor on the right is
$\gamma$. It maps $\Lambda(L,L')$ to
$\Lambda(\gamma L,\gamma'L')$.

First assume $\det(1-\gamma)\ne0$. For the action on the source
factor, $\det_{\mathbb C}(1-T_{\gamma,1}^{-1})
=\prod_\tau\tau(\det(1-\gamma))^b$. Since $a$ is even,
\[
 \det\gamma\,\overline{\det(1-\gamma)}=\det(1-\gamma).
\]
If $b=2r$ this identity gives
$|w^{2r}|/w^{2r}=\tau(\det\gamma)^{-r}$ for
$w=\tau(\det(1-\gamma))$. Equations (CA.15) and (CA.24) prove
\[
 P_{\gamma,1}\theta_0
   =\left(\prod_\tau\tau(\det\gamma)^{-b/2}\right)
       J_{A_0,T_{\gamma,1}A_0}\theta_0.
 \tag{CA.26}
\]
For the target factor use $\gamma'^{-1}=\bar\gamma'^{t}$.
The conjugate determinant identity has the opposite phase, yielding
\[
 P_{1,\gamma'}\theta_0
   =\left(\prod_\tau\tau(\det\gamma')^{a/2}\right)
       J_{A_0,T_{1,\gamma'}A_0}\theta_0.
 \tag{CA.27}
\]
If one rectangular dimension is zero, these are identities on
constant functions with empty products; the Gaussian argument is
unnecessary.

The restriction excluding eigenvalue one can be removed exactly.
There are infinitely many rational norm-one elements
$z=(1+t\sqrt{-d})/(1-t\sqrt{-d})$, $t\in\mathbb Q$.
Choose $z\ne1$ outside the finitely many inverse eigenvalues of
$\gamma$. Then $\gamma=(z\gamma)(z^{-1}1)$ and each factor has
$1-$factor invertible. Apply (CA.26) to each factor, use
(CA.17), then compose the two $J$ operators by CA4. Determinants
multiply, so (CA.26) holds for every $\gamma$. The same argument
applies to (CA.27). Combining the two commuting actions gives
\[
 P_{\gamma,\gamma'}\theta_0
    =\rho(\gamma,\gamma')J_{A_0,T_{\gamma,\gamma'}A_0}\theta_0,
 \quad
 \rho(\gamma,\gamma')=
       \prod_\tau\tau(\det\gamma)^{-b/2}
                       \tau(\det\gamma')^{a/2}.
 \tag{CA.28}
\]

\paragraph{CA9: all three equivariance laws for the actual normalized family.}
Let $\chi$ satisfy the literal \texttt{IsChi} laws. For self-dual
$L,L'$ define the source scalar and actual family
\[
 d_\chi(L,L')=
  \chi(\operatorname{detId}(\mathcal O_E^a,L))^{b/2}
  \chi(\operatorname{detId}(\mathcal O_E^b,L'))^{-a/2},\qquad
 e(L,L')=d_\chi(L,L')\,\texttt{thMoved}(L,L').
 \tag{CA.29}
\]
These determinant ideals have norm one. Multiplicativity and the
principal norm-one law imply that their character values are
nonzero, and
\[
 d_\chi(\gamma L,\gamma'L')=
       \rho(\gamma,\gamma')d_\chi(L,L').
 \tag{CA.30}
\]
Indeed the determinant ideal of an image is the old ideal times
$(\det\gamma)$, and \texttt{IsChi} assigns to this principal ideal
$\prod_\tau\tau(\det\gamma)^{-1}$. This determines both signs of
the exponents in (CA.30).

Using the actual equality (CA.16), naturality (CA.17), and (CA.28),
\[
\begin{split}
 P_{\gamma,\gamma'}e(L,L')
  &=d_\chi(L,L')J_{TA_0,T\Lambda(L,L')}
                        P_{\gamma,\gamma'}\theta_0\\
  &=\rho(\gamma,\gamma')d_\chi(L,L')
                        J_{A_0,T\Lambda(L,L')}\theta_0
   =e(\gamma L,\gamma'L').
\end{split}
 \tag{CA.31}
\]
For a change in just one lattice, (CA.13), (CA.15), and
$\operatorname{detId}(L,M)\operatorname{detId}(M,N)=\operatorname{detId}(L,N)$ similarly give
\[
\begin{split}
 e(L,L'')&=\chi(\operatorname{detId}(L',L''))^{-a/2}
            I_{\Lambda(L,L'),\Lambda(L,L'')}e(L,L'),\\
 e(M,L')&=\chi(\operatorname{detId}(L,M))^{b/2}
            I_{\Lambda(L,L'),\Lambda(M,L')}e(L,L').
\end{split}
 \tag{CA.32}
\]
Equations (CA.31)--(CA.32) are all three clauses of
\texttt{Base.EqP}, hence of the actual
\texttt{normalisedFamily\_satisfies\_equivarianceLaws} theorem.
The construction starts with the supplied Base's own
$\theta_0=\texttt{thZero}$, furnished by VI10. It neither chooses
a different family nor assumes the sought covariance of this one.

\paragraph{CA10: lift independence and covariance of the actual indicator.}
Assume the literal \texttt{LaHyp}, $\mathfrak A\ne0$, and $M>0$.
Thus $S$ contains no conjugate pair, the period ideal is a unit at
each $P\in S$, and $e(L,L')$ is fixed by the period lattice.
Use rational coordinates and put
$\mathcal P=\mathfrak c\mathfrak A\operatorname{Hom}(L,L')$.
For $c>0$ the two half lattices have the following local readings:
\[
\begin{array}{c|ccc}
 &P\in S&\bar P\in S&P,\bar P\notin S\\ \hline
 (A_c)_P&c^{-1}\mathcal P_P&c\mathcal P_P&\mathcal P_P\\
 (B_c)_P&c\mathcal P_P&c^{-1}\mathcal P_P&\mathcal P_P.
\end{array}
 \tag{CA.33}
\]
These localizations determine the global modules by LP1. Opposite
exponents cancel in the conjugate pairing, so both half lattices
are integral for the cocycle $c(x,y)$ of (CA.3). This follows
also by applying the integral pairing to $wx,y$, since $w$ is
integral and preserves each lattice.

Let $x\in M^{-1}\operatorname{Hom}_S(L,L')$. There exists
$t\in A_{Mc}$ with $t-x\in c\operatorname{Hom}_S(L,L')$.
To see this without an extra approximation hypothesis, clear a
common denominator for $x$ and an integral basis of the full
period module. Chinese remainders in this finitely generated
module give $y$ congruent to $(Mc)^{-1}x$ modulo $\mathcal P_P$
at the finitely many $P\in S$, and integral in $\mathcal P_Q$
at every other prime. Equivalently one may perform the CRT in
a common integral denominator module; it imposes only finitely
many residue conditions. Then $t=Mc\,y$ has the stated properties
by the three cases of (CA.33), using
$\mathcal P_P=\operatorname{Hom}(L_P,L'_P)$ on $S$.

Take any such $t$ and any representatives
$R\subset B_c$ modulo $B_c\cap\mathcal P$. Put
\[
 n_c=[\mathcal P\cap A_c:B_c\cap\mathcal P\cap A_c]>0.
\]
The literal finite Fourier definition reduces to
\[
 \operatorname{thInd}_{S,L,L',c,x}
    =\frac1{n_c}W_{-t}\sum_{y\in R}W_y e(L,L').
 \tag{CA.34}
\]
Indeed its inner representative set is
$(\mathcal P\cap A_c)/(B_c\cap\mathcal P\cap A_c)$.
The indicator of $B_c$ selects only its zero class. At that
class its Fourier phase is one, since $y\in B_c$ and the
cocycle is integral on $B_c$. Thus the Fourier coefficient is
$v/(n_c\sqrt v)$; the additional source factor $1/\sqrt v$
cancels $v>0$, giving exactly $1/n_c$.

The resulting formula is independent of the chosen lift and
representatives. The sum $F=\sum_{y\in R}W_y e$ is $B_c$-fixed:
translation permutes its quotient representatives and the period
and cocycle phases are one. Let $s$ be the source's original
lift in $A_c$. At $P\in S$, $t-s\in c\mathcal P_P$; at the
conjugate primes $s\in c\mathcal P_P$ and
$t\in Mc\mathcal P_P$; elsewhere both are integral.
Consequently $t\in A_c$ and $t-s\in B_c$.
The integral cocycle on $A_c$ gives
$W_{-t}W_{t-s}=W_{-s}$, so $W_{-t}F=W_{-s}F$.
This proves (CA.34) for the actual chosen-lift definition, not
just for a conveniently redefined indicator.

Now let $\gamma,\gamma'$ be rational unitary matrices.
The rational map $\phi(x)=\gamma'x\gamma^{-1}$ carries
$\mathcal P$, $A_c$, $B_c$, $M^{-1}\operatorname{Hom}_S$,
and $c\operatorname{Hom}_S$ to the corresponding modules for
$\gamma L,\gamma'L'$. This follows directly from the defining
mapping property and from localization; scalar denominators
commute with $\phi$. It preserves $n_c$, sends $t,R$ to valid
lift and representative data, and its complex map satisfies
$P_TW_y=W_{\phi(y)}P_T$. Applying CA9 to (CA.34) proves
\[
 \operatorname{thInd}_{S,\gamma L,\gamma'L',c,\phi(x)}(z)
       =\operatorname{thInd}_{S,L,L',c,x}(T^{-1}z).
 \tag{CA.35}
\]
This is exactly \texttt{ThIndEq}, with the source inverse on
the target factor. Equation (CA.30) is separately the literal
\texttt{ChiLaw}. Neither statement discards the determinant
factor before normalizing the jet.

\paragraph{CA11: the actual step has a finite algebraic jet field.}
Let $B$ be the given base, $X:\texttt{PreX}(B)$, and
$I:\texttt{ColIn}(B,X)$ satisfy its literal \texttt{Ok}.
Fix an index $a_0$, a word $u$, arbitrary steps $i,j$, and
$M>0,n\ge0$. Then
\[
 \operatorname{FinAlg}\!\left(
  \operatorname{jetSet}(I.A,N_i,N_j,I.S_b(a_0,u),M,n)\right).
 \tag{CA.36}
\]
The source also supplies \texttt{HasLetters}; we retain that
antecedent when matching \texttt{finAlg\_step}, although the
analytic proof below does not need an auxiliary stage.

First verify its actual period hypotheses. Every $N_i$ is even.
The selected set $S=I.S_b(a_0,u)$ is the union of the type-prime
set and the auxiliary/tuple set. A type prime and its conjugate
cannot both occur, because the two factors of its prescribed
norm element are coprime. The same applies to each auxiliary
factor. Tuple primes are chosen on only one side of a split
pair; primes of different supported rational types cannot collide.
This proves that $S$ has no conjugate pair. Their exact local
conditions say $q_P=\operatorname N P$ is prime and
$q_P\notin P^2$: for the prescribed norm elements this follows
by factoring their product and using the degree and norm count;
for tuple primes it is the literal split, degree-one, unramified
condition. These are the cases in \texttt{bad\_primes}.

At such a $P$ the relative different of $E/K$ is a unit.
In fact $q_P\notin P^2$ says that the ramification index over
$\mathbb Q_{q_P}$ is one; the residue extension is separable,
so the completed integral extension is unramified and its trace
pairing is perfect. The same holds over the intermediate field
$K$. The actual integral half-different $\mathfrak b$ satisfies
$\mathfrak b\bar{\mathfrak b}=\mathfrak D_{E/K}$.
Its two nonnegative local valuations must therefore be zero.
Thus $\mathfrak A=\mathfrak b^{-1}$ localizes to
$\mathcal O_{E,P}$. Since the curve ideal here is exactly
$\mathfrak c=\mathcal O_K$, this is the required
\texttt{AAUnit}; it is not inferred from self-duality alone.
VI2 supplies full period lattices, their self-duality, integral
quadratic multiplier, and finite generation. VI10 supplies the
actual fixed theta vector. Together these prove every
\texttt{LaHyp} clause, and $\mathfrak A\ne0$.

VI11, using only the independent finite presentation CA2, gives
\texttt{ESmooth} and \texttt{JetAlg}. CA9--CA10 give the actual
\texttt{VacEq}, \texttt{ThIndEq}, and \texttt{ChiLaw}.
The numbers $c_D(\tau)$ are algebraic: \texttt{AmbIn.Ok.dgen}
places each one in its specified number-field image.

Here is the finite-field argument, making the degree assertion
explicit. TL1 gives finitely many rational-unitary classes of
self-dual lattices in each rank. Choose representatives for
$L,L'$, finitely many positive divisors $c$ of $M$, and then
representatives for
\[
 M^{-1}\operatorname{Hom}_S(L,L')/
                    c\operatorname{Hom}_S(L,L').
 \tag{CA.37}
\]
The last quotient is finite by localization of a full lattice
and CRT; only the primes of $S$ dividing $Mc$ contribute.
Changing $x$ by its denominator module leaves the actual lift
predicate, hence the chosen indicator, unchanged.
For each of these finitely many functions retain its normalized
derivatives at zero of every total order at most $n$.
VI11 makes this a finite algebraic set. Adjoin also the finitely
many $c_D(\tau)$ and the images under every $\tau$ of a fixed
$\mathbb Q$-basis of $E$. The resulting field $K_0$ is generated
by finitely many algebraic numbers and contains every embedded
rational-unitary matrix entry and determinant.

For arbitrary $L,L',x$, choose unitary matrices carrying the
chosen lattice representatives to them. Equation (CA.35)
expresses their indicator by linear substitution in a selected
one. Repeated differentiation of a linear substitution is a
finite sum of derivatives of the same total order, multiplied
by products of the substitution's matrix entries. All these
entries lie in $K_0$. Equation (CA.30) shows that the ratio of
normalizing characters is in $K_0^\times$. Finally
\texttt{jetC} multiplies by the prescribed product of
$c_D(\tau)$ powers. Thus every normalized order-$n$ jet lies
in $K_0$, proving (CA.36). The polynomial here is the finite
Taylor component of a fixed degree, or its linear substitution;
no theta function is asserted to be a polynomial.

There is no hidden \texttt{LocStage}/\texttt{NStage} assumption
in this argument. The source proof of \texttt{finAlg\_step}
introduces arbitrary complete stages to access notation indexed
by $D7$. Its host-level \texttt{VacIn} already contains exactly
the period pair, half-point, rational isometry, complex reading,
and correction roots used above. It defines the same
\texttt{thMM}, \texttt{thZero}, \texttt{thMoved}, and
\texttt{eNat}; the source \texttt{toRec} and \texttt{*\_rec}
identities simply copy these fields. Every CA1--CA10 and VI
argument is expressed in those fields, the host, and its
embeddings. Applying those arguments directly to $I.J.V$ proves
(CA.36) for the literal $I.A=I.J.\texttt{amb}(B,I.\chi)$.
No unused complete-stage witness is required or claimed.

\paragraph{CA12: one finite generator set gives all canonical jets.}
Fix a preliminary datum $P$, character $\chi$, and finite twist
family $\Xi$. The actual finite bad-prime set is
$\mathcal B=\texttt{badP}(P,\Xi)$. For each word $u$, step
$i<m$, and subsets $S,R\subset\mathcal B$, put
\[
 q_R=\prod_{Q\in R}\operatorname NQ,\qquad
 M_{u,R}=q_R^{m+1}\bigl(T_uq_R^{m+1}\bigr).
 \tag{CA.38}
\]
The source \texttt{stepSet} uses this exact modulus, ranks
$N_i,N_{i+1}$, degree $\operatorname{totDeg}(P,u,i)$, and the
specific ambient datum $\texttt{amb0}(P.\mathrm{host})$.
Its \texttt{allSet} is the union only over those indices for
which \texttt{FinAlg(stepSet)} holds. There are finitely many
indices: the word and step sets are finite, and $S,R$ range over
two finite power sets. Taking the union of the finite algebraic
generator sets proves \texttt{FinAlg(allSet)}. This argument
does not silently assert finiteness for every arbitrarily chosen
prime set; the restriction in the actual definition is retained.

Let $Z=\texttt{gens}(P,\chi,\Xi)$ be the actual chosen finite
generator set. Every element is algebraic, and if a subfield
$K_0\subset\mathbb C$ contains $Z$, then it contains
\texttt{allSet}, by the defining universal property of field
generation. The fallback value $Z=\varnothing$ is never used
when $\Xi$ is finite.

Now let $I$ satisfy \texttt{Ok}, \texttt{IsTw}, and
$I.J=\texttt{amb0}(B.\mathrm{host})$, and assume
\[
 Z(B.\mathrm{toPre},I.C.\chi_\Psi,\operatorname{range}I.\xi)
                   \subset B.\mathrm{embLL}(B.LL).
 \tag{CA.39}
\]
Canonicality says that both exponents in $I.\mathrm{cst}$
are $m+1$, so its literal modulus is exactly (CA.38) for
$R=\texttt{rowSG}$; the type factor $T_u$ and all $q_Q$ are
positive. Its $S_b$ and $\texttt{rowSG}$ are subsets of
$\mathcal B$. Crucially, CA11 proves \texttt{FinAlg} for this
actual indexed step, so it satisfies the restriction in
\texttt{allSet}. The weight-to-degree identity of TL19 gives
$e.\mathrm{degree}=\operatorname{totDeg}(P,u,i)$ for every
multi-index tested by \texttt{JetsIn}. Therefore each such jet
lies in \texttt{allSet}, hence in the field image in (CA.39).
This proves the actual \texttt{jetsIn\_of\_isTw}.

\paragraph{CA13: the same-data canonical evaluation theorem.}
The preceding arguments close the analytic premises retained in
CI19/TL21. We spell out the quantifiers of the resulting
\texttt{EvLg} statement. Fix a source context $K$. Its
\texttt{large\_2} witnesses may be chosen as follows.
The fixed coefficient-field requirement contains a generator
of TL14's common frame-value field. The full coefficient-field
requirement contains the PA13 period-field generators, TL15's
finite roots and square roots for the prescribed finite twist
family, the auxiliary-prime square roots, and CA12's $Z$.
Take the union of these finite sets; all its entries are
algebraic. The extra finite prime and avoidance requirements
can be empty. These sets are chosen before the later base and
field-containment antecedents, exactly as \texttt{large\_2}
requires. The twist family is finite by TL21's support argument.

Suppose a base $B$, frame $f$, and twist datum $T$ lie over
$K$ and satisfy those field requirements. Retain the literal
frame-inclusion antecedent
\[
 i_X'(\operatorname{frame}_{X'}(a))
       =\operatorname{frame}_{X_m}(b)
       \quad\text{whenever the finite indices have equal values},
 \tag{CA.40}
\]
and suppose $\texttt{PreX}(B.\mathrm{toPre})$ is nonempty.
For every prescribed selection $P_v$ satisfying \texttt{IsSel},
the following construction uses this same $B,f,T,P_v$.

First orient the auxiliary factors of a supplied \texttt{PreX}
by TL16, choosing each factor or its conjugate according to
$f.P$. Replace only its type-prime elements by the context's
actual $\pi_T$. Since $c_u=\pi_T(u_0)/\pi_T(u)$, cancellation
of the nonzero denominator gives
$c_u\pi_T(u)=\pi_T(u_0)$; auxiliary orientation is unchanged.
Together with (CA.40), this is exactly both fields of
\texttt{Pins}. It asserts no new frame identity for arbitrary
raw frames.

Choose the actual $J=\texttt{amb0}(B.\mathrm{host})$;
PA9 proves its \texttt{AmbIn.Ok}. TL14--TL15 and the fixed/full
field containments supply the actual character values, tuples,
selected primes, and finite root data. CI15--CI17 construct
the local index records and all their nonvanishing and character
clauses, with initial $M=A=B=m$ and depth $2m$ off the auxiliary
primes. The fieldwise join of CI19 therefore gives a literal
$I_0:\texttt{ColIn}$ with
\[
 I_0.C=f,\quad I_0.\mathrm{sel}=P_v,\quad
 I_0.S=\{p\in\texttt{offSet}(B):p\ne p_\delta,p_\delta'\},
 \quad\operatorname{range}I_0.\xi=T.\mathrm{set}(f),\quad I_0.J=J.
 \tag{CA.41}
\]
The depth $2m$ satisfies the exact \texttt{Deep} bound needed
by CI18. Hence the specified replacement
$I=\texttt{canon}(I_0)$ satisfies \texttt{Ok} and
\texttt{IsTw}, preserving every field in (CA.41).
The same full-field requirement contains CA12's generators,
so CA12 supplies \texttt{JetsIn} for this $I$, with no change
of $J$, frame, or twist range. TL13 then also supplies
\texttt{Laws}, although the literal \texttt{evLg\_canonical}
conclusion does not require it as an extra conjunct.

TL14 supplies \texttt{ValLL0} for every later assembly
$C,W,V,S,N,G$ on this same base and frame; PA13 supplies the
same \texttt{FieldOK}. These are the final two conjuncts of
the source conclusion. A universal assertion over later
assemblies is not their existence assertion.

The construction works for every context $K$ already allowed
by the source. Thus the outer \texttt{Ev} wrapper uses
$\varepsilon_0=1$ and zero additional size thresholds, as in
\texttt{ev\_of\_forall}; all its original context and inner
field-containment quantifiers remain. This proves the actual
\texttt{evLg\_canonical} implication with its stated antecedents.
It neither creates an arbitrary preliminary datum, removes
(CA.40), nor constructs \texttt{Later}. The SE component
separately describes its stage suppliers and its still-retained
\texttt{LocStage}/\texttt{NStage} boundary. No inference to an
unconditional proof of RH or GRH is made here.
